% Generated draft: requires mathematical and source review.
\documentclass[11pt]{article}
\usepackage[margin=25mm]{geometry}
\usepackage{fontspec}
\setmainfont{Latin Modern Roman}
\usepackage{amsmath,amssymb,mathtools,unicode-math}
\setmathfont{Latin Modern Math}
\usepackage{enumitem,needspace,xurl,hyperref}
\setlength{\parindent}{0pt}
\setlength{\parskip}{5pt}
\setlength{\emergencystretch}{3em}
\newcommand{\N}{\mathbb{N}}
\newcommand{\Z}{\mathbb{Z}}
\newcommand{\Q}{\mathbb{Q}}
\newcommand{\R}{\mathbb{R}}
\newcommand{\C}{\mathbb{C}}
\newcommand{\F}{\mathbb{F}}
\newcommand{\A}{\mathbb{A}}
\newcommand{\PP}{\mathbb{P}}
\newcommand{\E}{\mathbb{E}}
\newcommand{\eps}{\varepsilon}
\newcommand{\dd}{\,\mathrm{d}}
\begin{document}
\section*{UnsolvedMath ID: 2308015}
\textbf{Source problem code:} catalogue entry
\subsection{Definitions and mathematical statement}
Let $\mathbb T=\partial\mathbb D$ with normalized arc measure, and let $H^2\subset L^2(\mathbb T)$ be the Hardy space, the closed span of $1,z,z^2,\ldots$. Write $H^2_-=L^2\ominus H^2$ and let $P_-$ be its orthogonal projection. In the scalar-symbol convention, a bounded symbol $\phi\in L^\infty(\mathbb T)$ gives the Hankel operator
\[
H_\phi:H^2\to H^2_-,\qquad H_\phi f=P_-(\phi f).
\]
In the bases $z^k$ and $z^{-j-1}$ its matrix is $\big(\widehat\phi(-j-k-1)\big)_{j,k\geq0}$, constant on antidiagonals. Equivalently one may study bounded Hankel matrices on $\ell^2(\mathbb Z_{\geq0})$, identified with $H^2$.

A bounded operator $H$ is trace class when it is compact and its singular values, the eigenvalues of $(H^*H)^{1/2}$ with multiplicity, satisfy
\[
\sum_{j\geq1}s_j(H)<\infty.
\]
\textbf{Problem (F. Holland).} Give a necessary-and-sufficient characterization of the scalar Hankel operators on $H^2$ that are trace class, in terms of their symbols or Hankel coefficients. Only the negative Fourier part of $\phi$ affects $H_\phi$; a symbol characterization must respect that nonuniqueness.

The supplied research summary reports this characterized by Peller through a Besov-type symbol condition and labels the problem solved with moderate confidence. It does not state the exact Besov space, norm, or theorem citation, so no more specific equivalence is reconstructed from that phrase alone.
\subsection{Short English statement}
Characterize which scalar Hardy-space Hankel operators have summable singular values. The supplied summary attributes the known Besov-space characterization to Peller without a precise theorem reference.
\subsection{Sources}
\begin{itemize}
\item Hayman, \emph{Research Problems in Function Theory} (2018), Problem 8.15, attributed to F. Holland, \url{https://arxiv.org/abs/1809.07200}; source-TeX extraction, accessed 2026-07-29. The 2018 update was retained for status triage.
\item Peller, trace-class Hankel characterization by a Besov-type symbol condition, as attributed in the research summary; no title, year, theorem number, or bibliographic identifier is supplied.
\item Local registry \texttt{data/unsolvedmath/problems.json}, permanent ID $2308015$, source code catalogue entry, \emph{Research Problems in Function Theory --- Problem 8.15}. The record labels this solved and SOLVED-IN-LITERATURE, with moderate-confidence research wording and NEEDS\_REVIEW source-status and rights qualifications.
\end{itemize}
\end{document}
